% ============================================
% LUC-13
% Extraído de main.pandoc.tex
% ============================================
\fragtitle{luc}{L}{LUC}{13}{A Noite do Diagnóstico}
\markboth{LUC-13 --- A Noite do Diagnóstico}{}

\noindent \textit{Documento extraído do arquivo pericial.}

\vspace{0.5em}\noindent\rule{\textwidth}{0.5pt}\vspace{0.5em}

\noindent \textbf{Abril de 2026 --- Belo Horizonte}\\
\noindent \textbf{Gravação de áudio --- arquivo L260-13A.wav}\\
\noindent \textbf{Duração:} 47 minutos e 32 segundos.\\
\noindent \textbf{Status:} não transcrito. A Dra. Lúcia Mendes não autorizou a transcrição.

\noindent \textit{Nota do arquivista: este arquivo de áudio foi encontrado no computador de Dra. Lúcia Mendes, datado de 12 de abril de 2026. É o primeiro registro pessoal dela sobre o caso. A gravação foi feita em casa, às 3:47 da manhã. Transcrevemos os primeiros 4 minutos e 12 segundos --- o restante é ininteligível.}
\noindent \textbf{[00:00:00] Lúcia:} Teste. Teste de áudio. L260... L260... treze. A. Gravando.

\noindent \textbf{[00:00:12]\textit{ É a terceira noite. Terceira noite que eu acordo no mesmo horário. 3:14. Sempre 3:14. Não adianta ajustar o relógio. Não adianta não olhar a hora. Eu }sei} que são 3:14 antes de olhar. Sinto na respiração. No peso do ar.

\noindent \textbf{[00:00:41]} Tentei todas as hipóteses. Ansiedade antecipatória. Condicionamento clássico. Sugestão pós-hipnótica induzida pela leitura. Tudo o que eu ensinaria num curso de psicopatologia para residentes. Quinze anos de clínica e eu sei exatamente como um paciente descreve um despertar noturno recorrente e como se diferencia um transtorno do sono de uma ideação persecutória.

\noindent \textbf{[00:01:18]} Mas não é um paciente. Sou eu.

\noindent \textbf{[00:01:25]\textit{ E o que eu sinto quando acordo não é ansiedade. É }expectativa}. Como se algo fosse acontecer. Como se algo estivesse \textit{se preparando} para acontecer. O corpo espera. O corpo \textit{sabe} de algo que a mente ainda não autorizou.

\noindent \textbf{[00:02:03]} O cheiro voltou hoje. Tinha passado o dia inteiro sem sentir. Pensei: passou. Foi sugestão. O cérebro cansado fabrica estímulos onde não existem. Aí desliguei a luz do quarto, deitei, fechei os olhos... e o cheiro estava no travesseiro. Doce. Não doce de perfume. Doce de matéria vegetal murchando dentro de água parada. Doce de buquê esquecido num copo sujo até a água ficar preta.

\noindent \textbf{[00:02:47]} Levantei, troquei a fronha. Cheirei o travesseiro novo. Limpo. Lavanda. O cheiro continuou. Não estava no travesseiro. Estava dentro do meu nariz.

\noindent \textbf{[00:03:12]} Acendi todas as luzes. Abri a janela. Respirei fundo. O ar da rua estava frio, normal, urbano. Poluição. Diesel. Asfalto molhado. Coisas reais. Fechei a janela. O cheiro voltou na mesma intensidade.

\noindent \textbf{[00:03:41]} Ele está \textit{em mim}. Não no ambiente. Na mucosa. No fundo da garganta. No lugar onde a memória olfativa se aloja e não sai nunca mais. Como o cheiro da casa da infância. Como o cheiro de alguém que se foi.

\noindent \textbf{[00:04:12]} Fui ao banheiro e olhei no espelho. O olho direito...

\noindent \textit{Aqui a gravação se interrompe por 23 segundos. O que se segue foi recuperado parcialmente por restauração de áudio. A voz de Lúcia está diferente: mais baixa, mais lenta, como se falar exigisse esforço físico.}
\noindent \textbf{[00:04:35] Lúcia:} ...está mais claro. Não muito. O suficiente. A íris perdeu pigmento na borda externa. Ficou cor de... não sei. O amarelo ainda não chegou. É um tom que não tem nome. Um tom que eu não quero nomear.

\noindent \textbf{[00:05:11]} Toquei o olho com a ponta do dedo. A córnea estava morna. Olhos não são mornos. Olhos são frios. A temperatura da superfície ocular é 32 graus. O que eu senti era mais. Bem mais.

\noindent \textbf{[00:05:44]\textit{ Fiz um diário. Não. Fiz um }registro}. Uma planilha. Sintomas, horários, intensidade, frequência. Durou três dias. No terceiro dia, percebi que a planilha estava igual à tabela que Joaquim descreve no diário dele. A mesma estrutura. As mesmas colunas. As mesmas palavras. A minha letra, mas a mesma estrutura.

\noindent \textbf{[00:06:28]\textit{ Eu não copiei. Eu juro que não copiei. A tabela }veio}. Veio pronta. A mão escreveu sozinha.

\noindent \textbf{[00:07:01]} Rasguei a planilha. Joguei no lixo. Depois tirei do lixo, queimei na pia, lavei as cinzas. Não quero que ninguém encontre.

\noindent \textbf{[00:07:33]} Se eu lesse um relato como este vindo de uma paciente, meu diagnóstico seria: ideação paranóide com somatização, possível quadro dissociativo leve, estresse pós-traumático deflagrado por contato com material traumático.

\noindent \textbf{[00:08:12]} Mas o que eu sinto não é paranoia. Paranoia é achar que alguém está te observando. O que eu sinto é diferente. Eu sinto que alguém \textit{já está dentro}. O que Joaquim descreve como \textquotedblleft{}o peso atrás do esterno\textquotedblright{}, \textquotedblleft{}a respiração que não é minha\textquotedblright{}, \textquotedblleft{}a boca que se move sem eu mandar\textquotedblright{}. Eu sei o que ele quis dizer. Porque estou sentindo.

\noindent \textbf{[00:08:52]} Na madrugada eu sentei pra escrever. Uma coisa que eu queria registrar. Peguei o caderno. Abri na primeira folha. A mão escreveu: \textit{tem tanta fome aqui, Lúcia. Tanta fome que eu nem sei onde caber.}

\noindent \textbf{[00:09:07]\textit{ Essa frase está no diário. É a primeira fala da entidade. Eu }li} ela. Mas quando a mão escreveu, a letra não era minha. A letra era a mesma descrição que o Dr. Oliveira fez na carta dele: inclinada pra esquerda, \textquotedblleft{}a\textquotedblright{} aberto, pressão irregular. A letra de Joaquim.

\noindent \textbf{[00:09:41]\textit{ Eu não copiei. Eu não }precisava} copiar. A mão sabia o caminho sozinha. A mão lembrou de um padrão que eu vi no papel e reproduziu sem autorização. A mão tem memória própria.

\noindent \textbf{[00:10:03]} O mesmo padrão. A mesma estrutura. Eu escrevendo a frase de Joaquim. Eu sentindo o que ele sentiu. Eu vendo o que ele viu. Eu.

\noindent \textbf{[00:10:28]} Antes eu falava \textquotedblleft{}Joaquim descreve\textquotedblright{}, \textquotedblleft{}Joaquim relata\textquotedblright{}, \textquotedblleft{}o paciente diz\textquotedblright{}. Agora eu falo \textquotedblleft{}a vala\textquotedblright{}, \textquotedblleft{}o olho\textquotedblright{}, \textquotedblleft{}ela\textquotedblright{}. O nome próprio dele está sumindo do meu vocabulário. Ele não é mais o paciente. Ele é a testemunha. A testemunha que depõe dentro de mim.

\noindent \textbf{[00:11:12]} Daqui a pouco eu vou falar \textquotedblleft{}nós\textquotedblright{}.

\noindent \textit{[A gravação registra um ruído longo, semelhante a folhas sendo viradas depressa.]}
\noindent \textbf{[00:12:04]} Não vem. Não vem. Não vem. Ela não pode estar aqui. Ela não pode estar dentro de mim. Eu sou psiquiatra. Eu tenho 42 anos. Eu não acredito em possessão. Eu não acredito em contágio. Eu não acredito em nada, em nada, em nada...\margnota{o som de\\folhas sendo\\viradas.\\Quem está\\virando as\\páginas?}

\noindent A partir deste ponto, a gravação se deteriora. A voz de Lúcia alterna entre silêncios longos e frases curtas, desconectadas. A qualidade do áudio piora --- há um ruído de fundo que os técnicos descrevem como \textquotedblleft{\textit{semelhante a vento passando por um duto estreito\textquotedblright{}.}

\noindent \textbf{[00:12:04] Lúcia:} O diário está no escritório. Eu tranquei a porta.

\noindent \textbf{[00:12:31]} Não adianta. Sinto ele do outro lado da parede. A capa. O peso. A lombada.

\noindent \textbf{[00:13:02]\textit{ Ele }irradia}.

\noindent \textbf{[00:13:28]} Uma hora dessas... 3:47 da manhã... e eu sei que ele está acordado. O diário está acordado. As páginas estão se movendo sozinhas no escuro. A tinta está escorrendo para lugares que eu não escrevi.

\noindent \textbf{[00:14:15]} Isto não é científico. Isto não é clínico. Isto é o relato de uma mulher adulta, com pós-doutorado, cujo cérebro está mentindo para ela.

\noindent \textbf{[00:14:52]} \textit{(longa pausa)} E se não for mentira?

\noindent Os 33 minutos seguintes são compostos principalmente de silêncio, respiração irregular e ruído de fundo. Aos 47 minutos e 12 segundos, a gravação termina com o que parece ser a voz de Lúcia sussurrando: \textquotedblleft{\textit{Paciente 1.262.\textquotedblright{} A gravação foi encerrada manualmente em 47:32. Nenhum outro arquivo de áudio foi encontrado no computador.}

\noindent Nota do arquivista: na manhã seguinte a esta gravação, Lúcia Mendes ligou para o Arquivo Público Mineiro e solicitou a devolução do diário. Disse que precisava \textquotedblleft{\textit{confirmar uma informação}\textquotedblright{}} na contracapa. O diário foi retirado por ela pessoalmente no mesmo dia. Ela não o devolveu até 13 de maio.}

\vspace{0.5em}
\noindent\textit{\textbf{↪ Rotas:} \rota{DOC-01} \rotasep \rota{LUC-02} \rotasep \rota{EPI-01} \rotasep \rota{LUC-14}}

}